\documentclass[11pt,fleqn]{book}
\usepackage[english]{babel}
\usepackage{lipsum}
\usepackage{xcolor}
\usepackage{tikz}
\usepackage{mathtools,amsfonts,amssymb,amsthm}
\usepackage[most]{tcolorbox}
\definecolor{ocre}{RGB}{52,177,201}
\definecolor{ultramarine}{RGB}{0,45,97}
\definecolor{mybluei}{RGB}{0,173,239}
\definecolor{myblueii}{RGB}{63,200,244}
\definecolor{myblueiii}{RGB}{199,234,253}


\usepackage{blindtext}
\usepackage{fourier}
\usepackage[explicit,calcwidth]{titlesec}

\renewcommand\thechapter{\arabic{chapter}}

\newcommand\chapnumfont{%
  \fontsize{38}{13}\color{myblueii}\selectfont%380-130
}

\newcommand\chapnamefont{%
  \normalfont\color{white}\scshape\small\bfseries
}

\titleformat{\chapter}
  {\normalfont\huge\filleft}
  {}
  {0pt}
  {\stepcounter{chapshift}%
  \begin{tikzpicture}[remember picture,overlay]
  \fill[myblueiii]
    (current page.north west) rectangle ([yshift=-2cm]current page.north east); %svetlomodra
  \node[
      fill=mybluei,
      text width=2\paperwidth,
      rounded corners=1cm,
      text depth=2cm,%temnomodra
      anchor=center,
      inner sep=0pt] at (current page.north east) (chaptop)
    {\thechapter};%
  \node[
      anchor=south east,
      inner sep=5pt,%notranji odmik
      outer sep=0pt] (chapnum) at ([xshift=-20pt]chaptop.south)
    {\chapnumfont\thechapter};
  \node[
      anchor=south,
      inner sep=0pt] (chapname) at ([yshift=2pt]chapnum.south)
  {};
  \node[
      anchor=north east,
      align=right,
      inner xsep=0pt] at ([yshift=0.1cm]chapname.east|-chapnum.south) %polozar NASLOVA
  {\parbox{\textwidth}{\raggedleft#1}};%tole poglej da bo cel naslov not
  \end{tikzpicture}%
  }

\titleformat{name=\chapter,numberless}
  {\normalfont\huge\filleft}
  {}
  {0pt}
  {\begin{tikzpicture}[remember picture,overlay]
  \fill[myblueiii]
    (current page.north west) rectangle ([yshift=-13cm]current page.north east);
  \node[%BRISI
      fill=mybluei,
      text width=2\paperwidth,
      rounded corners=6cm,
      text depth=18cm,
      anchor=center,
      inner sep=0pt] at (current page.north east) (chaptop)
    {};%
  \node[ %BRISI
      anchor=south east,minimum width=2in,
      inner sep=0pt,
      outer sep=0pt] (chapnum) at ([xshift=-20pt]chaptop.south)
    {};
  \node[ %BRISI
      anchor=south,
      inner sep=0pt] (chapname) at ([yshift=2pt]chapnum.south)
  {};
  \node[ %BRISI 
      anchor=north east,
      align=right,
      inner xsep=0pt] at ([yshift=-0.5cm]chapname.east|-chapnum.south)
  {\parbox{\textwidth}{\raggedleft#1}};
  \end{tikzpicture}%
  }

\titlespacing*{\chapter}{0pt}{0pt}{0in} %odmik texta
\titlespacing*{name=\chapter,numberless}{0pt}{0pt}{0in}

\titleformat{\section}
  {\addtolength{\titlewidth}{4pc}\normalfont\Large\sffamily\bfseries}
  {\llap{\colorbox{myblueii}{\parbox{2cm}{\strut\color{white}\hfill\thesection}}\hspace{20pt}}}
  {0em}{#1}
  [{\titleline*[l]{\hspace*{\dimexpr-2cm-20pt-2\fboxsep\relax}\color{myblueii}\titlerule[1.5pt]}}]
\titleformat{\subsection}
  {\addtolength{\titlewidth}{2pc}\normalfont\large\sffamily}
  {\llap{\colorbox{myblueii}{\parbox{2cm}{\strut\color{white}\hfill\thesubsection}}\hspace{20pt}}}
  {0em}{#1}
  [{\titleline*[l]{\hspace*{\dimexpr-2cm-20pt-2\fboxsep\relax}\color{myblueii}\titlerule[1.5pt]}}]

\usetikzlibrary{calc}
\pagestyle{plain}

\newcounter{chapshift}
\addtocounter{chapshift}{-1}

\newcommand\BoxColor{myblueii}

\def\subsectiontitle{}
\renewcommand{\sectionmark}[1]{\markright{\sffamily\normalsize#1}{}}
\renewcommand{\subsectionmark}[1]{\def\subsectiontitle{#1}}


\usepackage{etoolbox,fancyhdr}

\pagestyle{fancy}

\renewcommand{\headrule}{{\color{myblueii}%
\hrule width\headwidth height\headrulewidth depth\headrulewidth}}

\renewcommand{\chaptermark}[1]{\markboth{\sffamily\normalsize\bfseries \ #1}{}}
\renewcommand{\sectionmark}[1]{\markright{\sffamily\normalsize#1}{}}
\fancyhf{} \fancyhead[LE,RO]{\sffamily\normalsize\colorbox{myblueii}{\color{white}\sffamily\bfseries\strut\quad\thepage\quad}}
\fancyhead[LO]{\textcolor{mybluei} \rightmark%

%%%%%%%%%%%%%%%%%%%%%%%%CODE FOR SIDE BAR%%%%%%%%%%%%%%%%%%%%%%%%%

\begin{tikzpicture}[overlay,remember picture]
  \node[fill=\BoxColor,inner sep=-1pt,rectangle,text width=1cm,
    text height=28cm,align=center,anchor=north east]
  at ($ (current page.north east) $)
  {\rotatebox{90}{\parbox{28cm}{%
\centering\textcolor{black}{\bfseries\scshape\thechapter.\leftmark \hspace{2cm}\color{myblueiii} company name }\hspace{2cm}\rightmark}}};
% 	   \raggedright\textcolor{black}{\bfseries\scshape\thechapter.\leftmark \hspace{2cm}\rightmark}\hspace{2cm}\color{myblueiii} company name}}};
  \end{tikzpicture}}
\fancyhead[RE]{\textcolor{mybluei}\leftmark%
  \begin{tikzpicture}[overlay,remember picture]
  \node[fill=\BoxColor,inner sep=0pt,rectangle,text width=1cm,
    text height=28cm,align=center,anchor=north west]
  at ($ (current page.north west) $)
  {\rotatebox{90}{\parbox{28cm}{%
\centering\textcolor{black}{\bfseries\scshape\thechapter.\leftmark \hspace{2cm}\color{myblueiii} company name}\hspace{2cm}\rightmark}}};
  \end{tikzpicture}}


%%%%%%%%%%%%%%%%%%%%%%%% END CODE FOR SIDE BAR%%%%%%%%%%%%%%%%%%%%%%%%%

\renewcommand{\headrulewidth}{.5pt}
\addtolength{\headheight}{2.5pt}
\newcommand{\footrulecolor}[1]{\patchcmd{\footrule}{\hrule}{\color{#1}\hrule}{}{}}
\renewcommand{\headrulewidth}{.5pt}
\addtolength{\headheight}{2.5pt}


\fancypagestyle{plain}{%
  \fancyhf{}%
  \renewcommand{\headrulewidth}{0pt}
  \renewcommand{\footrulewidth}{0pt}
}

\makeatletter
\renewcommand{\cleardoublepage}{
\clearpage\ifodd\c@page\else
\hbox{}
\vspace*{\fill}
\thispagestyle{empty}
\newpage
\fi}
\makeatother

\begin{document}
	

	
\chapter{\textcolor{green}{Uvod}}
\section{\textcolor{green}{Zakonodaja}}
\section{\textcolor{green}{Zasnova sistema za pripravo bazenske vode}}

\chapter{Klorirna postaja}
\section{Jama za nevtralizacijo}
\section{Menjava klornih jeklenk}
\section{Opis naprave za doziranje klora} %\section{Dezinfekcija s klorom (kloriranje)}  urska je tuki naporej ime gruntam ce bi pustu kloriranje 
\subsection{Način delovanja}
\subsection{Priključitev injektorja na vod pogonske vode}

\chapter{\textcolor{red}{Klor in reakcija z vodo}} % kr opis

\chapter{Avtomatski vakuumski klorinator - AKN}%General+Complete device description
\section{Tipična uporaba sistema}
\section{Seznam delov ter opis delovanja}

\chapter{Razvoj elektronike}
\section{Mikrokrmilnik - Kinetis KL05 MKL05Z32VFM4}
\section{Senzor tlaka}
\subsection{Piezoelektrilni senzor tlaka}
\subsection{Ničelna napetost}
\subsection{Senzor tlaka - City sensors OC18-104}
\section{Procesorska ploščica} %pr tem delu elektronsko opises zakaj je ksn element na svojem mestu in kaj dosezemo z tem (kako vezje deluje in zakaj)
\subsection{Napajanje/glajenje vhodne napetosti}
xc,nbcxykljbcnxynmbbmbcxmy
\subsection{Nastavitev merilnega območja senzorja (R3/R2)}
\subsection{Invertirajoči ojačevalnik (preračunavanje meritev)}
\subsection{Izhodni signal za kontrolno enoto}
\subsection{Programiranje mikrokrmilinika}
\section{Vhodno-izhodna ploščica} %pr tem delu elektronsko opises zakaj je ksn element na svojem mestu in kaj dosezemo z tem (kako vezje deluje in zakaj)
\subsection{Matrični zaslon (5 $\times$ 7) - HCMS-3903} %kakšen zaslon mamo in kako deluje,...
\subsection{Nastavitev merilnega območja (Tipki HI/LO)}
\subsection{Vizualna predstavitev odčitkov} % ledice na kerem obmocju svetjo,..
\section{\textcolor{blue}{CAD program za risanje tiskanih vezij Eagle}}
\section{Tiskanina ter montažna shema}
\section{Seznam materiala}

\chapter{Razvoj mehanike} 
\blindtext [5]
\section{Prednji del ohišja}
\blindtext [5]
\section{Notranji del ohišja}
\blindtext [5]
\blindtext [5]
\section{Zunanji del ohišja}

\chapter{Razvoj programske opreme} 
\section{Opis programa} %flow chart opises
\section{Opis funkcij} % kako naprava dejansko deluje/kako se jo prizge
\section{Vizualizacija meritev}


\chapter{Testiranje izdelke/meritve}

\chapter{Zaključek}
\includegraphics[width=\textwidth]{example-image} 
\blindtext[5]
\chapter{CAD program za risanje tiskanih vezij Eagle}

\end{document}
